\chapter{Összegzés} % Conclusion
\label{ch:sum}

Összességében elmondható, hogy egy olyan játékszoftvert sikerült fejlesztenem, amely egyaránt szórakoztató és tanulságos.

Felhasználói szempontból a játékélmény változatos és élvezetes, ugyanis a játékos mindig újabb kihívások elé kerül, akár egy, vagy többjátékos módról legyen szó. A látványos háromdimenziós vizuális megjelenítés által a felhasználó tényleg egy igazi labirintusban érezheti magát, még ha nem is élethű a grafika. Emellett szabadságot ad az egyéni kreativitás kifejezéséhez a beépített pályaszerkesztő és a konfigurálható labirintusgenerátor.

A fejlesztő pedig megismerheti egy jól strukturált, szabadon bővíthető játékmotor működését, bővítheti a szakmai ismereteit. Belátást nyer OpenGL és az SDL világába, valamint a többszálú programozás és a hálózati kommunikáció működésébe. Gyakorlatban alkalmazva kaphat betekintést egy labirintusgenerátor algoritmus C++ implementációjába és tetszése szerint módosíthat rajta. Emellett számtalan lehetőséget kap arra is, hogy kifejezze a kreativitását a játékmotor továbbfejlesztésével és más, akár hálózatos játékot fejlesszen, ami nem feltétlenül labirintusjáték.

A szoftvert az egyetemen szerzett tudást és szemléletmódot alkalmazva valósítottam meg, a fejlesztés közben sikerült új ismereteket, valamint tapasztalatot szereznem. Számomra izgalmas feladatokkal szolgált és sikerélményt nyújtott egy-egy főbb rész elkészülése. Úgy gondolom, hogy sikerült bővítenem a játékfejlesztői tudásomat, amit majd később alkalmazni tudok az iparban. Remélem, hogy a szakdolgozatommal másoknak segítséget és motivációt adhatok a játékfejlesztésben.